\begin{table}[H]
\centering
\caption{Categories of External Covariates (adapted from Tuominen et al. \cite{tuominen_forecasting_2024})}
\label{tab:covariates}

% 1. Use \begin{tabularx}{\linewidth}{...} to make the table fit the page.
% 2. Use {lX} for the columns:
%    'l' = a standard, left-aligned column for the short category names.
%    'X' = a special, auto-wrapping column for the long description.
\begin{tabularx}{\linewidth}{lX} 
\toprule
% 3. Separate your headers with an ampersand (&)
\textbf{Category} & \textbf{Description} \\
\midrule
Hospital beds & Hourly bed availability data from 24 healthcare facilities in the catchment area. \\
Traffic & Hourly vehicle counts from 33 bidirectional observation stations in the Pirkanmaa region. \\
Weather & Ten historical weather variables from the nearest meteorological observation station. \\
Public events & A time series counting the number of daily public events organized in the Tampere area. \\
Website visits & Website traffic for two hospital domains and Google search trend data for "Acuta". \\
Calendar variables & Temporal features such as day of the week, month, national holidays, and holiday lags. \\
Technical analysis & Features engineered from the target variable itself, such as moving averages. \\
\bottomrule
\end{tabularx}
\end{table}